\documentclass[11pt,a4paper]{article}
\usepackage[review]{acl}
\usepackage{times,latexsym}
\usepackage{booktabs}

\title{ARR Reproducibility Checklist\\(Supplementary Material)}
\author{Anonymous ACL submission}

\begin{document}
\maketitle

\section*{Responsible NLP Research Checklist}

\subsection*{A. Limitations and Risks}

\begin{enumerate}
\item \textbf{Did you discuss the limitations of your work?} \\
\textit{Yes.} A dedicated Limitations section discusses: single-dataset scope (with DebugBench cross-validation), four models across three architecture families (1.5B--13B), seed-dependent variance on DeepSeek, the invalidated full\_trace run, generator--evaluator overlap, self-familiarity bias, error distribution assumptions, and task scope limited to verdict prediction.

\item \textbf{Did you discuss any potential risks of your work?} \\
\textit{Yes.} The Ethics Statement notes that the work uses public benchmarks, involves no human subjects, and has minimal environmental impact. Verdict prediction is a developer-facing tool with no direct end-user impact.
\end{enumerate}

\subsection*{B. Did you use or create scientific artifacts?}

\begin{enumerate}
\item \textbf{Did you cite the creators of artifacts you used?} \\
\textit{Yes.} HumanEval \cite{chen2021evaluating}, MBPP-sanitized \cite{austin2021program}, DebugBench \cite{tian2024debugbench}, Qwen-2.5-Coder-7B/1.5B-Instruct \cite{qwen2025coder}, CodeLlama-13B-Instruct \cite{roziere2024code}, DeepSeek-Coder-6.7B-Instruct \cite{deepseek2024coder}, LoRA \cite{hu2022lora}, and TRL SFTTrainer are all cited.

\item \textbf{Did you discuss the license or terms for use and/or distribution of any artifacts?} \\
\textit{Partial.} All base models and datasets are publicly available. HumanEval is released under MIT license; MBPP under CC-BY-4.0; DebugBench under MIT. Model licenses: Qwen-2.5-Coder under Apache 2.0; DeepSeek-Coder under the DeepSeek License.

\item \textbf{Did you discuss if your use of existing artifacts was consistent with their intended use?} \\
\textit{Yes.} HumanEval and MBPP are standard code generation benchmarks; we use their canonical test cases for verdict evaluation, consistent with their evaluation purpose. DebugBench is designed for debugging evaluation.

\item \textbf{Did you discuss the steps taken to check whether the data contains any information that names or uniquely identifies individual people or offensive content?} \\
\textit{N/A.} All datasets contain only programming problems and solutions with no personal information.
\end{enumerate}

\subsection*{C. Computational Experiments}

\begin{enumerate}
\item \textbf{Did you report the number of parameters in the models used, the total computational budget, and computing infrastructure?} \\
\textit{Yes.}
\begin{itemize}
\item Models: Qwen-2.5-Coder-7B/1.5B-Instruct (7B/1.5B parameters), CodeLlama-13B-Instruct (13B parameters), DeepSeek-Coder-6.7B-Instruct (6.7B parameters). Only LoRA adapters are trained (rank=16, $\alpha$=32, targeting all linear layers).
\item Hardware: 2$\times$ RTX PRO 6000 Blackwell (${\sim}$96\,GB VRAM each); each experiment uses a single GPU.
\item Total compute: ${\sim}$320 GPU-hours across all conditions, models, seeds, and iterative refinements (${\sim}$170h for successful experiments; individual runs 0.5--2.2h). See Appendix~B for per-condition wall-clock times.
\end{itemize}

\item \textbf{Did you report all hyperparameters?} \\
\textit{Yes.} Appendix~B (Table~5) reports: LoRA rank (16), alpha (32), dropout (0.05), target modules (all linear layers), epochs (3), per-device batch size (4), gradient accumulation steps (8), effective batch size (32), learning rate ($2 \times 10^{-4}$), LR scheduler (linear decay), warmup ratio (0.1), weight decay (0.01), max sequence length (2,048 tokens).

\item \textbf{Did you report descriptive statistics about your results?} \\
\textit{Yes.} All accuracy values include Wilson 95\% confidence intervals. Pairwise comparisons use paired bootstrap resampling ($B{=}10{,}000$, seed=42) with reported $p$-values. We report multi-seed results (seeds 42, 123, 456) for the exception\_only condition; other conditions report single-seed results. Standard deviations are provided where multiple seeds are available.

\item \textbf{If you used existing packages, did you report the implementation, model, and their versions?} \\
\textit{Yes.} HuggingFace TRL SFTTrainer with LoRA via PEFT. Base model identifiers: \texttt{Qwen/Qwen2.5-Coder-7B-Instruct} and \texttt{deepseek-ai/deepseek-coder-6.7b-instruct}.
\end{enumerate}

\subsection*{D. Human Annotations}

\begin{enumerate}
\item \textbf{Did you use human annotators?} \\
\textit{No.} All labels (pass/fail) are derived from automated test execution. No human annotation was involved.
\end{enumerate}

\subsection*{E. Data and Code Availability}

\begin{itemize}
\item \textbf{Code}: Training, evaluation, and trace preprocessing scripts will be released upon acceptance.
\item \textbf{Data}: Execution traces, VOTA-compressed traces, and train/test splits will be released. Raw traces are generated deterministically from publicly available benchmarks using Python \texttt{sys.settrace}.
\item \textbf{Models}: LoRA adapters for all conditions will be released on HuggingFace.
\end{itemize}

\subsection*{F. Reproducibility Summary}

\begin{table}[h]
\centering
\small
\begin{tabular}{lc}
\toprule
\textbf{Item} & \textbf{Status} \\
\midrule
All hyperparameters reported & Yes \\
Random seeds specified & Yes (42, 123, 456) \\
Error bars / CIs reported & Yes (Wilson 95\% CI) \\
Statistical tests reported & Yes (bootstrap, McNemar) \\
Computing infrastructure reported & Yes \\
Training time reported & Yes \\
Code release planned & Yes \\
Data release planned & Yes \\
Model release planned & Yes \\
\bottomrule
\end{tabular}
\caption{Reproducibility checklist summary.}
\end{table}

\bibliography{references}

\end{document}
